% Copyright 2025 by Lernik Asserian
% This file may be distributed and/or modified.

\documentclass[11pt,letterpaper,reqno]{amsart}
\usepackage{amsmath}
\usepackage{amssymb}
\usepackage{amsthm}
\usepackage{booktabs}
\usepackage{graphicx}
\usepackage{xspace}
\usepackage[dvipsnames]{xcolor}
\usepackage{setspace}
\usepackage[letterpaper,margin=1in,headheight=15pt]{geometry}
\usepackage{bm}
\usepackage{hyperref}
\definecolor{hyperlinkblue}{RGB}{0,0,238}
\hypersetup{colorlinks=true,urlcolor=hyperlinkblue,linkcolor=hyperlinkblue,citecolor=hyperlinkblue}
\allowdisplaybreaks

% Makes the subsection go onto a new line
% \makeatletter
% \def\subsection{\@startsection{subsection}{2}%
%   \z@{.5\linespacing\@plus.7\linespacing}{.5\linespacing}%
%   {\normalfont\bfseries}}
% \makeatother

% Customizing the inclusion of codes in Appendix
\usepackage{listings}
\lstdefinestyle{mystyle}{
    backgroundcolor=\color{lightgray!20}, % A light gray background
    commentstyle=\color{green!60!black},
    keywordstyle=\color{blue},
    numberstyle=\tiny\color{gray},
    stringstyle=\color{purple},
    basicstyle=\fontfamily{pcr}\selectfont\footnotesize, % Use a courier-like font
    breakatwhitespace=false,         
    breaklines=true,                 
    captionpos=b,                    
    keepspaces=true,                 
    numbers=left,                    
    numbersep=5pt,                  
    showspaces=false,                
    showstringspaces=false,
    showtabs=false,                  
    tabsize=2
}
\lstset{style=mystyle}

\numberwithin{equation}{subsection}
\newtheorem{thm}{Theorem}
\newtheorem{lem}[thm]{Lemma}
\newtheorem{prop}[thm]{Proposition}
\newtheorem{cor}[thm]{Corollary}
\theoremstyle{definition}
\newtheorem{Example}[thm]{Example}
\newtheorem{Definition}[thm]{Definition}
\numberwithin{thm}{subsection}
\newtheorem{problem}{Problem}

% labeling of figures and tables
\newcounter{figtab}[subsection]
\setcounter{figtab}{0}
\makeatletter
\@addtoreset{figure}{subsection}
\@addtoreset{table}{subsection}
\makeatother
\renewcommand{\thefigure}{\thesubsection.\arabic{figtab}}
\renewcommand{\thetable}{\thesubsection.\arabic{figtab}}
\let\oldfigure\figure
\def\figure[#1]{\oldfigure[#1]\stepcounter{figtab}}
\let\oldtable\table
\def\table[#1]{\oldtable[#1]\stepcounter{figtab}}

\usepackage{biblatex}
\addbibresource{references.bib}


\title{How Do Acoustic Features Relate to the Emotions We Hear in Music?}
\author{Caden Che and Kevin Zuniga}
\thanks{Exploratory analysis prepared for MATH 195, Stanford University.}
\begin{document}
\onehalfspacing
\begin{abstract}
We investigate how measurable properties of recorded music relate to its emotional character, using valence and arousal as separate responses. Our shared interest in energetic songs that sound sad motivated an initial focus on electronic dance music. Following a review of the candidate recordings, we broadened the question to music across genres. This midterm report presents the revised framework alongside the completed pilot analysis of 92 provisional dance-style candidates from MERGE. Seven acoustic features supported a modest improvement in arousal prediction using ridge regression, while valence prediction did not improve over a training-mean baseline. These results concern the pilot sample; the broader sample has not yet been evaluated. Ordinary least squares and ridge regression provide our starting models, with source separation and an LSTM proposed as extensions. Two selected songs will later connect the fitted predictions to our listening observations.
\end{abstract}
\maketitle
\input{pilot/01_introduction}
\input{pilot/02_background_and_literature_review}
\section{Methodology and Technical Implementation}
\input{pilot/03_1_mathematical_model}
\input{pilot/03_2_data_selection_audit}
\input{pilot/03_3_implementation}
\input{pilot/03_4_results}
\input{pilot/04_conclusion}
\printbibliography
\appendix
\section{Supplementary Descriptive Figures}
\begin{figure}[ht]
\centering
\includegraphics[width=0.75\linewidth]{figures/05_descriptive_correlations.pdf}
\caption{Descriptive Pearson correlations between the seven acoustic features and supplied labels. These correlations do not adjust for sample composition or artist-credit groups.}
\end{figure}
\begin{figure}[ht]
\centering
\includegraphics[width=\linewidth]{figures/06_split_sensitivity.pdf}
\caption{Changes in error relative to each split's training-mean baseline. Positive values favor ridge. The ten evaluations reuse the same observations.}
\end{figure}
\end{document}

